% !TEX program = xelatex
% AI-effects 프로젝트 작업 노트 --- 산업별 LCE(labor cost equivalent) 연구를 위한 한국 자료와 준비사항.
% 계획 노트: 실제 자료 수집·분석은 아직 하지 않음. 자료의 공개 수준과 공표 시기는 신청 전 재확인 필요.
% 관련: paper 19(한국 대화량), paper 21(SOC–KSCO 연계표), paper 23(KSCO 결합과 특이사항)
% 컴파일: xelatex (kotex 필요)
\documentclass[11pt]{article}

\usepackage{fontspec}
\usepackage{kotex}
\usepackage[a4paper,margin=2.2cm]{geometry}
\usepackage{longtable}
\usepackage{booktabs}
\usepackage{array}
\usepackage{xcolor}
\usepackage{hyperref}
\usepackage{enumitem}

\hypersetup{colorlinks=true,linkcolor=blue!50!black,urlcolor=blue!50!black}
\newcolumntype{L}[1]{>{\raggedright\arraybackslash}p{#1}}
\setlength{\LTleft}{0pt}
\setlength{\LTright}{0pt}
\renewcommand{\arraystretch}{1.15}
\setlength{\tabcolsep}{4pt}

\title{산업별 LCE 연구를 위한 자료와 준비사항\\[2mm]
\large --- 필요한 한국 통계, 분류 연계, LCE 설계, 분석상 유의점 ---}
\author{AI-effects 프로젝트 (작업 노트)}
\date{\today}

\begin{document}
\maketitle

\begin{abstract}
\noindent 한국 직업별 AI 대화량과 지표(paper 23)를 산업별 직업 분포로 산업에 배분하고, 산업별 임금과 결합해 산업별 LCE(labor cost equivalent)를 계산할 계획이다. 이어 산업별 LCE와 고용, 임금, 노동시간, 부가가치, 노동생산성(노동시간 기준)의 관계를 분석한다. 본 노트는 이 연구에 필요한 한국 통계, 분류 연계, LCE 계산에서 정할 사항, 분석상 유의점, 저장소 준비사항을 정리한 계획 노트다. 실제 자료 수집과 분석은 아직 하지 않았다. 자료의 공개 수준(마이크로데이터의 직업$\cdot$산업 코드 자릿수 등)과 공표 시기는 작성자의 지식에 근거한 것이므로, 신청 전에 MDIS와 각 기관의 공표 일정에서 다시 확인해야 한다.
\end{abstract}

\tableofcontents

% ============================================================
\section{연구 흐름}
% ============================================================

\begin{enumerate}[leftmargin=*]
  \item \textbf{보유 자료}: 한국 직업별 업무용 AI 대화량과 AI 지표(월 $\times$ KSCO 8차 세분류, \texttt{kor\_ksco8\_conv\_ai\_indicators}; SOC 단위 \texttt{kor\_soc6\_conv\_ai\_indicators}).
  \item \textbf{배분}: 산업 $\times$ 직업 고용(또는 근로시간) 분포로 직업별 대화량을 산업에 나눈다.
  \item \textbf{가격 부여}: 산업 $\times$ 직업별 대화량 $\times$ 대화당 절감 시간에 산업 $\times$ 직업별 시간당 임금(노동비용)을 곱해 산업별 LCE를 만든다.
  \item \textbf{결합과 분석}: 산업 분류를 연계해 고용, 임금, 노동시간, 부가가치, 노동생산성(시간당 부가가치)과 결합하고 관계를 분석한다.
\end{enumerate}

% ============================================================
\section{필요한 한국 통계}
% ============================================================

\subsection{산업 $\times$ 직업 고용 분포(배분 가중치)}

\begin{longtable}{L{3.6cm}L{1.9cm}L{5.2cm}L{5.3cm}}
\caption{산업 $\times$ 직업 고용 분포 자료}\label{tab:emp}\\
\toprule
자료 & 기관 & 장점 & 한계 \\
\midrule
\endfirsthead
\toprule
자료 & 기관 & 장점 & 한계 \\
\midrule
\endhead
\textbf{지역별고용조사}(반기, MDIS) & 통계청 & 표본이 커서 산업 $\times$ 직업 교차가 안정적. 자영업자 포함. 주당 취업시간 있음 & 공개용 자료의 직업$\cdot$산업 코드 자릿수 확인 필요(세분류는 원격접근일 수 있음) \\
경제활동인구조사(월) & 통계청 & 월별 시계열 & 공개 분류가 거침. 교차 분포용으로는 표본 부족 \\
\textbf{고용형태별근로실태조사}(MDIS) & 고용노동부 & 근로자 단위로 산업 $\times$ 직종 $\times$ 임금 $\times$ 근로시간이 함께 있음 & 임금근로자만 포함(자영업자 제외) \\
인구주택총조사 20\% 표본 & 통계청 & 가장 세밀한 산업 $\times$ 직업 분포 & 5년 주기(최근 2020, 다음 2025) \\
고용보험 DB & 고용정보원 & KECO 직종 $\times$ 산업, 행정자료 전수. KECO 연계는 이미 준비됨 & 접근 제한. 피보험자만 포함 \\
\bottomrule
\end{longtable}

\noindent 권장: 기본 배분은 지역별고용조사, 임금 분석에는 고용형태별근로실태조사를 쓴다. 인구총조사는 검증용으로 둔다.

\subsection{산업 $\times$ 직업별 임금$\cdot$노동비용}

\begin{itemize}[leftmargin=*]
  \item \textbf{고용형태별근로실태조사}: 시간당 임금(정액$\cdot$초과급여, 특별급여 포함 총액)을 산업 $\times$ 직종 단위로 구할 수 있다. LCE의 기본 단가로 쓴다.
  \item \textbf{기업체노동비용조사}(고용노동부): 직접 노동비용(임금)과 간접 노동비용(퇴직급여, 사회보험료, 교육훈련비 등)의 비율을 산업별로 준다. 임금을 ``노동비용''으로 바꾸는 배율로 쓴다. 상용 10인 이상 회사법인만 대상이라는 점에 유의한다.
  \item \textbf{국민계정 경제활동별 피용자보수}(한국은행): 산업별 총 노동비용의 기준 총량이다. LCE를 산업 노동비용 대비 비율로 나타낼 때 분모로 쓴다.
  \item 자영업자 소득은 임금 자료에 없으므로, 노동소득분배율 조정 같은 별도 처리가 필요하다.
\end{itemize}

\subsection{산업별 성과 변수}

\begin{longtable}{L{3.4cm}L{6.2cm}L{6.4cm}}
\caption{산업별 성과 변수 자료}\label{tab:outcome}\\
\toprule
변수 & 자료 & 비고 \\
\midrule
\endfirsthead
\toprule
변수 & 자료 & 비고 \\
\midrule
\endhead
부가가치(명목$\cdot$실질) & 한국은행 국민계정 경제활동별 GDP & 국민계정 경제활동분류 기준. 연간 확정까지 시차가 김 \\
부가가치(세밀 산업) & 산업연관표(기본부문, 통합 중분류$\cdot$대분류) & 부속 고용표(취업자$\cdot$피용자 수)도 활용 가능 \\
사업체 단위 부가가치$\cdot$종사자$\cdot$급여 & 광업$\cdot$제조업조사, 서비스업조사, 경제총조사 & KSIC 세세분류까지 가능 \\
고용 & 지역별고용조사, 경제활동인구조사, 사업체노동력조사 & 가구 기준인지 사업체 기준인지 통일 필요 \\
임금$\cdot$근로시간 & \textbf{사업체노동력조사}(월), 고용형태별근로실태조사 & 월별 산업 중분류 시계열은 사업체노동력조사가 적합 \\
노동생산성(시간 기준) & 한국생산성본부 노동생산성지수, KIP DB & 노동투입 = 근로자 수 $\times$ 근로시간. 직접 계산할 때는 부가가치 $\div$ (취업자 $\times$ 연간 근로시간) \\
\bottomrule
\end{longtable}

\subsection{선택: 산업별 AI 이용률}

고용 비례 배분은 ``같은 직업이면 산업과 관계없이 AI를 같은 강도로 쓴다''는 가정이다. 산업별 AI 활용 정도를 보여 주는 자료로 이 가정을 점검하거나 보정할 수 있다. 후보는 통계청 기업활동조사의 AI 활용 항목, 과기정통부$\cdot$NIA의 인터넷이용실태조사와 생성형 AI 이용 조사, 한국노동연구원 등의 근로자 AI 이용 설문이다.

% ============================================================
\section{분류 연계 준비}
% ============================================================

\begin{longtable}{L{3.8cm}L{2.0cm}L{10.2cm}}
\caption{필요한 분류 연계}\label{tab:cw}\\
\toprule
연계 & 상태 & 필요 작업 \\
\midrule
\endfirsthead
\toprule
연계 & 상태 & 필요 작업 \\
\midrule
\endhead
SOC 2018 → KSCO 8차 세분류, KECO 2025 & \textbf{완료} & paper 21, 23 \\
\textbf{KSCO 7차 ↔ 8차} & 미확보 & KSCO 8차는 2025년 시행이라, 2024년 이전 지역별고용조사 등은 7차 코드일 가능성이 크다. 신구 연계표가 필요하다(\texttt{kostat\_ksco\_isco\_crosswalk/\_download\_manifest.md}에 따르면 사용자 폴더 \texttt{D:/data/AI/data}에 신구 연계표가 있음). \\
KSCO 세분류 → 소분류 집계 & 미작성 & 마이크로데이터가 소분류(3자리)까지만 공개될 경우 필요하다. 대화량은 합산하고, 지표는 고용 가중평균한다. \\
KSIC 10차 ↔ 11차 & 미확보 & 2024년 개정 반영 여부를 자료별로 확인한다. \\
\textbf{KSIC ↔ 국민계정 경제활동분류 / 산업연관표 부문} & 미확보 & 부가가치와 결합하는 데 꼭 필요하다(한국은행 연계표). \\
KECO 2018 ↔ 2025 & 미확보 & 과거 고용보험 DB를 쓰는 경우에만 필요하다. \\
\bottomrule
\end{longtable}

\noindent 공통 분석 단위는 국민계정 기준 산업 중분류 수준(약 30~60개)이 현실적이다. 산업 수가 적다는 점은 분석 설계의 제약이 된다(5절).

% ============================================================
\section{LCE 계산 설계에서 정할 사항}
% ============================================================

\begin{enumerate}[leftmargin=*]
  \item \textbf{정의(안)}: 산업 $i$, KSCO 직업 $k$에 대해
  \[
    \text{LCE}_i = \sum_k \text{대화량}_k \times \frac{\text{고용}_{ik}}{\text{고용}_k} \times \text{대화당 절감 시간}_k \times \text{시간당 노동비용}_{ik}.
  \]
  대화량 $\times$ 절감 시간의 곱은 SOC 단계에서 계산한 뒤 배분해야 한다(paper 23, 4.7절).
  \item \textbf{배분 기준}: 고용 인원과 근로시간 중 하나를 고르고, 기준 연도를 고정한다(예: 2023~2024년 평균). 결과 변수와 같은 해의 고용 분포를 쓰면 기계적 상관이 생긴다.
  \item \textbf{임금 단가}: 산업 $\times$ 직업별 임금을 쓰면 LCE와 산업 임금이 기계적으로 연동된다. 임금을 결과 변수로 볼 때는 직업별 전국 임금(산업 공통)으로 계산한 LCE를 함께 만든다.
  \item \textbf{대화량 관련 가정}:
  \begin{itemize}
    \item 대화당 메시지 수 $k$(1.5~10)에 따라 총량이 크게 달라지므로 민감도 분석이 필요하다.
    \item Claude.ai의 직업 구성을 모든 AI 제품에 적용했다는 가정이 들어 있다(paper 19).
    \item 시간 절감은 한국 값이 아니라 전 세계 값이다.
    \item 대화량은 2026년 4$\cdot$5월 두 달뿐이므로, 연간으로 환산하는 규칙을 정해야 한다.
  \end{itemize}
  \item \textbf{연계 특이사항 처리}: 대화량 미공개 KSCO는 결측(0이 아님)이고, 값이 같은 KSCO 묶음이 201개 있다(paper 23). 이를 산업 단계에서 어떻게 다룰지 정한다.
  \item \textbf{척도}: LCE를 절대액, 산업 노동비용 대비 비율, 취업자 1인당 값 중 무엇으로 쓸지 정한다. 산업 간 비교에는 비율 지표가 적합하다.
\end{enumerate}

% ============================================================
\section{분석 설계 시 유의점}
% ============================================================

\begin{itemize}[leftmargin=*]
  \item \textbf{시점 문제}: 대화량은 2026년 자료이지만, 2026년 연간 부가가치$\cdot$생산성(국민계정)은 2027년 이후에 나온다. 단기 분석은 월$\cdot$분기 자료(사업체노동력조사, 경제활동인구조사)로 하고, 연간 성과와의 관계는 사전 추세(2019~2025년)와의 단면 비교로 시작하는 것이 현실적이다.
  \item \textbf{역인과}: 생산성이 높은 산업이 AI를 더 많이 쓸 수 있다. 사전적 노출도(observed exposure, Eloundou)와 실제 사용(LCE)을 비교하거나, 노출도를 도구변수로 쓰는 방안을 검토한다.
  \item \textbf{표본 크기}: 산업 수가 30~60개라 검정력이 낮다. 산업 $\times$ 지역(지역별고용조사)이나 산업 $\times$ 직업 셀로 단위를 넓히는 방안을 함께 검토한다.
  \item \textbf{고용 개념 통일}: 가구조사(취업자, 자영업 포함)와 사업체조사(종사자, 임금근로자)를 섞지 않는다.
\end{itemize}

% ============================================================
\section{저장소 준비사항}
% ============================================================

\begin{itemize}[leftmargin=*]
  \item \textbf{원자료 폴더}: \texttt{data/raw/macro/} 아래에 자료별 하위 폴더를 만들고(예: \texttt{kostat\_lfs\_regional/}, \texttt{moel\_wage\_structure/}, \texttt{bok\_national\_accounts/}, \texttt{kpc\_productivity/}), \texttt{data/proc/macro/}에도 같은 구조를 둔다. 다운로드 기록(manifest)도 함께 작성한다. 폴더 이름은 새 결정이므로 작업 전에 확정한다.
  \item \textbf{신청$\cdot$접근}: MDIS 원격접근(세분류가 필요한 경우)이나 고용보험 DB 이용 신청에는 기간이 걸리므로 일찍 시작한다.
  \item \textbf{do 파일 구성(안)}:
  \begin{itemize}
    \item \texttt{01\_import\_*}: 자료별 읽기
    \item \texttt{01\_import\_*\_crosswalk}: KSCO 7↔8, KSIC↔국민계정
    \item \texttt{03\_merge\_industry\_occupation\_lce.do}
    \item \texttt{04\_analysis\_lce\_*}
  \end{itemize}
  \item \textbf{선행 노트}: LCE 정의, 배분 가정, 기준 연도, 산업 분류 수준을 확정하는 연구 노트를 먼저 작성한다.
\end{itemize}

\end{document}
